% =====================================================================
\chapter{Objections and Limits}\label{ch:objections}
% =====================================================================

This chapter performs a self-audit: it removes claims that cannot bear load from the
structure, and records explicitly what has been removed.

\section{List of Non-Claims}

The following list is binding: no other part of the paper may be read as supporting the
claims below.

\begin{enumerate}[label=(N\arabic*)]
\item The paper does \emph{not} claim that all large software eventually converges on one
      and the same high-level architecture (Hypothesis \ref{hyp:strong} is explicitly
      rejected).
\item The paper does \emph{not} claim that code volume by itself causes internal
      infrastructuralisation; it claims that the cause is the product of entity count and
      resource count, together with the recurrence of problems and the frequency of
      evolution.
\item The paper does \emph{not} claim that the six functions must be \emph{built
      in-house}; it claims that they must be \emph{defined}, and may be adopted from a host.
\item The paper does \emph{not} claim that Simon's near-decomposability is a universal law.
      The paper recasts it as a conditional, testable sufficient-statistic criterion
      (Definition \ref{def:neardecomp}) and concedes that the original source provides no
      general theorem about verification cost.
\item The paper does \emph{not} claim that minimum description length can decide what
      purpose a system ought to achieve. MDL is a model-selection principle, not a theory
      of intent.
\item The paper does \emph{not} claim that ``comprehension is compression'' is a theorem in
      every domain. Within inductive inference it is a defensible theoretical position;
      outside that range it is a strong heuristic.
\item The paper does \emph{not} claim that a mechanical system never errs. It claims that
      mechanical decision does not depend on the generator's subjective will, and that a
      mechanical system can still be defective, can fail, and can embody an incomplete
      model.
\item The paper does \emph{not} claim that trust can be fully mechanised. The \TCB{} can be
      reduced, verified and isolated, but it cannot be eliminated outright.
\item The paper does \emph{not} claim that the intent layer can be closed mechanically. It
      can be \emph{approached} by adversarial consensus, and the reliability of that
      approximation is governed by the correlated error mode.
\item The paper does \emph{not} claim that ``adding more agents to review'' raises
      reliability. By Corollary \ref{cor:diversity}, reliability rises only when the
      correlated error mode falls.
\item The paper does \emph{not} claim that layering can accommodate everything. A
      cross-cutting concern cannot be localised structurally into any single layer
      (Principle \ref{pr:crosscutting}).
\item The paper does \emph{not} claim that formal verification will become cheap
      \emph{automatically} in the AI era. That is a hypothesis worth investigating, not an
      established fact (\S\ref{sec:asymmetry}).
\item The paper does \emph{not} claim that compiler correctness equals system correctness.
      Linking, loading, the runtime, external libraries, hardware and the operating system
      need not lie within the scope of the proof.
\item The paper does \emph{not} claim that resource guarantees can be unified into a single
      method. Time complexity, symbolic resource bounds, worst-case execution time,
      statistical quantiles and direct measurement are evidence of different grades.
\item The paper does \emph{not} claim that ``no violation detected'' equals ``every
      execution satisfies''. Three states must be distinguished (Principle
      \ref{pr:three-states}).
\item The paper does \emph{not} claim to determine who is authorised to adopt a specification.
      Principle \ref{pr:authority} establishes that authorisation is a failure axis independent of
      reliability; which body holds the standing is an institutional fact that the framework
      locates but does not settle.
\item The paper does \emph{not} claim a model of how specification fidelity decays over time.
      Section \ref{sec:spec-binding} argues that the quantity must be time-indexed and that a
      re-confirmation rule is forced by Proposition \ref{prop:three-grounds}(i); the form of
      $\varepsilon_{\text{cons}}(\cdot)$ is left open.
\item The paper does \emph{not} claim that the design-effect expression \eqref{eq:design-effect} is a
      theorem about majority error. It is the standard approximation for a mean under exchangeable
      correlation, and is used here to make Corollary \ref{cor:diversity} quantitative, not to
      establish a bound.
\end{enumerate}

\section{Arguments Removed from the Load-Bearing Structure}

The paper deliberately rejects several arguments that are common in discussions of this
kind but do not hold. They are rejected not because they are useless, but because they
state heuristics as theorems, and thereby make the whole framework depend on unproved
conjectures.

\begin{longtable}{L{3.4cm}L{5.2cm}L{5.6cm}}
\caption{Arguments removed from the load-bearing structure}\label{tab:overreach}\\
\toprule
\textbf{Rejected argument} & \textbf{Why it fails} & \textbf{Replacement in this paper} \\
\midrule
\endfirsthead
\toprule
\textbf{Rejected argument} & \textbf{Why it fails} & \textbf{Replacement in this paper} \\
\midrule
\endhead
\bottomrule
\endlastfoot
Deriving from G\"odel's incompleteness theorem that ``a permanent gap between intent and
  specification necessarily exists'' &
  The incompleteness theorem concerns limits of provability in particular formal systems;
  it is not a theorem about the relation between human mental intent and formal
  specification &
  Replace it with three epistemic grounds: requirements are generative; specification is
  lossy compression; and the distinction between verification and validation
  (Proposition \ref{prop:three-grounds}) \\
Deriving from Rice's theorem that ``intent cannot be formalised'' &
  Rice's theorem concerns the undecidability of non-trivial semantic properties of
  programs; it has nothing to do with human intent &
  Use it only as a boundary on general synthesis and general decision algorithms (Q6 of
  Remark \ref{rmk:q1q8}) \\
Calling ``grading is easier than answering'' a dividend of $\mathrm{P}\neq\mathrm{NP}$ &
  Whether $\mathrm{P}$ equals $\mathrm{NP}$ is still open; the argument depends on an
  unproved conjecture &
  Use only the structural asymmetry implied by the definition of the class $\mathrm{NP}$
  (Proposition \ref{prop:asym-struct}) \\
Asserting as a mathematical theorem that ``the verifier must be minimal, deterministic and
  never isomorphic'' &
  A small trusted kernel is a valuable architectural principle; small does not mean free of
  defects &
  \TCB{} minimisation: make the untrusted part inspectable, and concentrate the unavoidable
  trust on a basis that is as small, and as explicit in its properties, as possible \\
Asserting that ``a dead thing does not betray'', hence mechanical assurance is absolutely
  reliable &
  A mechanical system can be defective, can fail, and can embody an incomplete model &
  Distinguish ``does not betray'' (does not depend on the generator's will) from ``cannot
  err'' (impossible) \\
Asserting that Simon's law of levels proves that every complex system necessarily organises
  itself into levels &
  The original source offers an explanation, not a universal law, and no general theorem
  about verification cost &
  Recast as a conditional, testable sufficient-statistic criterion
  (Definition \ref{def:neardecomp}) \\
Asserting that linear logic makes wasted resources structurally impossible &
  The exponential structure of linear logic still permits particular forms of duplication
  and discarding; linear use is not optimal use &
  Use a composite structure: usage constraints $+$ resource-bound analysis $+$ performance
  models $+$ runtime instrumentation $+$ qualified worst-case execution time analysis \\
Asserting that formal verification must become cheap in the AI era &
  Historical cost magnitudes show that specification design, model selection and proof
  architecture remain real work (Example \ref{ex:sel4-cost}) &
  Downgrade it to a hypothesis and give a measurable cost metric (criterion S5 of Chapter
  \ref{ch:program}) \\
Asserting that ``comprehension is compression'' is a universal law &
  Its rigorous formalisation applies to inductive inference, and beyond that range it
  degenerates into a heuristic &
  Restrict the scope of application and tie it to task-relative sufficiency
  (Definition \ref{def:sufficient}) \\
Asserting a substantive connection between this framework and the Langlands programme &
  That programme studies the functorial duality between arithmetic objects and automorphic
  representations; the ``free generation / deterministic admission'' this framework needs
  is only the \emph{smell of a duality}, not a functor. To write it in would substitute
  rhetoric for structure &
  Adopt a duality that pays its way: the LP/SMT duality. Farkas certificates, unsatisfiable
  cores and unrealisability certificates give a theorem-level correspondence between a
  synthesis problem and its infeasibility proof, and are already used in engineering \\
Asserting that the Landauer principle constrains software design &
  The scale of software is some twelve orders of magnitude away from that bound; the
  erasure bound is not a constraint surface for architecture &
  Adopt the memory wall and locality: moving one bit costs about $10^{9}$ times the erasure
  bound in energy, and this is the physics that actually fixes the memory hierarchy, cache
  policy and data layout \\
Asserting that homotopy type theory is a necessary foundation for this framework &
  That theory extends the foundations of mathematics, but gives no software guarantee that
  resource analysis or verified compilation does not already give &
  Restrict the role of type theory to what has been proved: dependent types expressing
  postconditions and invariants \cite{hott2013} \\
\bottomrule
\end{longtable}

\section{Possible Objections and Responses}

\paragraph{Objection 1: read-sets and oblivious sets cannot be measured precisely in real
systems, so the formalisation is decorative.}
Response: the read-set of Definition \ref{def:readset} can be counted directly from version
history, using ``the set of modules touched by each change'' as a proxy; it is a computable
quantity (Hypothesis $H_4$ exists precisely to test the validity of that proxy). The
oblivious set is not required here to be extracted automatically; it is required to be
\emph{explicitly declared}. It is a design artefact, not an analysis artefact. Once
declared, it can be checked in part by the leak-pair diagnostic (Corollary
\ref{cor:diagnose}).

\paragraph{Objection 2: the core conclusion that ``a large system grows a kernel'' is too
broad and almost unfalsifiable.}
Response: Theorem \ref{th:six} has explicit premises (more than one concurrently live
entity, finite contended resources, reachable faults), and its conclusion is falsifiable: a
well-defined system that satisfies the premises but lacks one of the six functions suffices
to refute it. The empirical part rests on Hypothesis $H_3$, whose falsifier is that ``lines
of code explain the data no worse than $n \times m$''.

\paragraph{Objection 3: the two-layer trust structure places too much responsibility on
consensus, and consensus often fails in practice.}
Response: this is exactly the problem Proposition \ref{th:correlated} addresses. The
proposition does not defend consensus; it gives the form of its failure and its design
variables. It also explains a common phenomenon: a ``review panel'' composed of models of
common origin is structurally not a guarantee, however large it is.

\paragraph{Objection 4: does treating the level as the unit into which verification
obligations decompose merely rename the existing practice of ``layered verification''?}
Response: in part, yes, and the paper concedes it. The difference lies in three points:
(a) the assignment rule given here (Definition \ref{def:obligation-assign}) is an
operational decision, not a slogan; (b) the paper explicitly introduces the class of
``obligations this layer cannot express'' and requires that it be assigned to consensus or
to human approval rather than ignored; (c) the paper gives a quantitative bound on the
reliability of the consensus layer, which makes ``assign it to consensus'' a decision that
can be assessed rather than merely trusted.

\paragraph{Objection 5: defining technical debt as a description-length gap is
uncomputable and therefore useless.}
Response: the value of the definition lies not in numerical computation but in the
direction and the criterion it gives (Principle \ref{pr:debt}(c)): deciding whether a
refactoring is worthwhile is equivalent to deciding whether it makes the system easier to
generate from a shorter description. That is a question that can be used in review, even
though $\DL^{\star}$ is uncomputable.

\paragraph{Objection 6: the paper mixes engineering, mathematics and social epistemology;
does it exceed the defensible range of a single discipline?}
Response: the criticism holds. Two mitigations are in place. First, every class of claim is
labelled with its strength (the strength trichotomy of Chapter \ref{ch:program}); second,
the citations that are \emph{used} are explicitly separated from those that are
\emph{rejected} (Table \ref{tab:relations}). Cross-disciplinarity is not itself a defect,
provided that no discipline's conclusion is used outside its scope.

\section{Open Problems}

The following questions are open:

\begin{enumerate}[label=(O\arabic*)]
\item \term{Independent measurability of $\Dom$ and $\Conv$}. For Conjecture
      \ref{conj:strength} to become a testable proposition, both quantities must be
      measurable independently of the conclusion. The paper gives only the direction, not a
      reproducible annotation protocol.
\item \term{Formalisation of the completeness argument for the six functions}. The proof of
      Theorem \ref{th:six} is a structural enumeration, not a derivation in a formal system.
      Writing it as a completeness theorem for a given operational-semantics framework
      remains unfinished work.
\item \term{Whether cross-layer semantics always compose}. Proposition \ref{prop:soundness}
      gives sufficient conditions, but whether those conditions can be satisfied
      simultaneously in systems containing nondeterminism, real-time constraints and
      reflexivity is not known.
\item \term{How to measure $q$ for the consensus layer}\label{op:q}. Proposition \ref{th:correlated}
      attributes reliability to the correlated error mode, but gives no operational method
      for estimating $q$. This is the crux of whether the direction can be put into
      practice. Section \ref{sec:consensus} supplies a companion quantity, the mean pairwise
      error correlation $\bar\rho$, which is estimable from data the framework already collects;
      $q$ remains the limit that is not. Appendix \ref{app:protocol} makes the same quantities
      addressable at the level of procedure---attestor roster, evidence provenance and dissent
      record are the fields from which $\bar\rho$ can be estimated and against which selection
      capture becomes visible---but it neither measures $q$ nor claims to.
\item \term{That the consensus protocol of Appendix \ref{app:protocol} works}. The appendix
      specifies a procedure and states its failure modes; it has not been built, run or
      evaluated, and no parameter in it was measured. Its role is to make the intent layer's
      obligations concrete enough to criticise, not to discharge them.
\item \term{Whether a joint semantics of resources and function is decidable}. A costed
      trace semantics can express joint constraints, but the size of its decidable fragment
      and the tightness of its bounds are unsettled.
\item \term{Compatibility of artefact binding with dynamically generated implementations}.
      If an implementation can be regenerated at runtime in response to load, the static
      identity assertion of Equation \eqref{eq:binding} must be replaced by a stronger
      question: which changes preserve the existing contract, which require re-verification,
      and how to fall back to a known safe state when verification fails.
\item \term{Consistency of level assignment}. Whether Definition
      \ref{def:obligation-assign} yields consistent assignments across different analysts
      requires empirical testing.
\end{enumerate}

\section{Methodological Limitations}

\begin{enumerate}[label=(L\arabic*)]
\item For every reference cited, the depth of checking is limited to the formal publication
      page, the author's or project's home page, and preprint metadata; the full text was
      \emph{not read through in full for each entry}. Every concrete number given in the
      text should therefore be treated as a \emph{second-hand report awaiting re-checking};
      several entries with recent nominal years are retained as recorded in the original
      source, and no independent checking was done here.
      A verification round was carried out for the first revision. The five empirical
      citations supporting Hypothesis $H_7$ were checked against the primary abstract pages:
      the four figures quoted (189 Lean tasks; $72.6\%$, $52.3\%$ and $4.9\%$ for the best
      model) and the 245 operating-system tasks in the second benchmark were all found
      \emph{in the abstracts themselves}, not merely in secondary summaries. Two venue
      corrections were made as a result: the traceable-refinement work is a workshop poster
      rather than a main-conference paper, and the mutation-driven method is an unrefereed
      preprint. Newly added citations were checked to the same depth, and the two entries
      whose full text was not read are marked ``metadata only'' in the bibliography.
      What the round did \emph{not} establish: no number was re-derived from a paper's tables
      or figures, and no full text was read end to end, so the grade of the numeric claims is
      still ``quoted from a primary abstract'', one step above the original second-hand
      report and short of independent replication.
\item The formalisation in this paper consists of \emph{working definitions established for
      the purposes of this paper}, and should not be mistaken for an existing unified
      standard. In particular, the three items ``oblivious set'', ``level contract
      four-tuple'' and ``level assignment of verification obligations'' are working
      concepts proposed here.
\item The three epistemic grounds (Proposition \ref{prop:three-grounds}) are arguments
      rather than theorems; the challengeable point is whether the three are exhaustive.
\item No experiments were performed. Every hypothesis in Chapter \ref{ch:program} remains
      in the ``to be tested'' state.
\end{enumerate}

\begin{principle}[Self-audit conclusion]\label{pr:self-audit}
After this scrutiny, three minimal independently testable propositions remain in the
load-bearing structure:
\begin{enumerate}[label=(\roman*)]
\item a structural asymmetry between verification and generation, implied by the definition
      of the class $\mathrm{NP}$;
\item the level is the unit into which assurance obligations decompose, because every
      obligation belongs to some observable domain;
\item the gap between intent and specification has epistemic grounds, and therefore cannot
      be closed mechanically in general.
\end{enumerate}
Every remaining claim (including ``compression'', ``emergence'' and ``internal
infrastructuralisation'') has either been downgraded to a falsifiable hypothesis or
confined to an explicit scope of application.
\end{principle}
